\documentclass[11pt,a4paper]{article}

\usepackage[utf8]{inputenc}
\usepackage[T1]{fontenc}
\usepackage{amsmath,amssymb}
\usepackage{booktabs}
\usepackage{graphicx}
\usepackage{xcolor}
\usepackage[margin=25mm]{geometry}
\usepackage{hyperref}
\hypersetup{colorlinks=true,linkcolor=blue!50!black,urlcolor=blue!50!black}

\graphicspath{{data/figures/}}
\newcommand{\tables}{data/tables}

\newcommand{\derived}{\textcolor{blue!50!black}{\textbf{[DERIVED]}}}
\newcommand{\measured}{\textcolor{green!35!black}{\textbf{[MEASURED]}}}
\newcommand{\opn}{\textcolor{orange!80!black}{\textbf{[OPEN]}}}

\title{The Fourier amplification symbol of the semi-Lagrangian quadratic fit\\
       \large The measured growth on hexahedral meshes is a boundary effect}
\author{Tomislav Mari\'c}
\date{2026-10-13}

\begin{document}
\maketitle

\begin{abstract}
The frozen-velocity power iteration measures a spectral radius above one on every mesh,
hexahedral meshes included: $\rho = 1.00441$ at Courant number $0.0164$, flat to $0.5\,\%$
over an eight-fold range of cell count. That invariance is the signature of a
constant-coefficient operator. This report derives the operator's Fourier symbol in closed
form. The interior operator on a uniform hexahedral mesh has $\rho = 1$ \emph{exactly}, for
every Courant number up to $0.95$, in one, two and three dimensions, and for both weight
anchors --- ninety-six configurations, no exception. The derivation is validated three ways:
symbolically against the known one-dimensional result, numerically against an independent
power iteration, and against the solver's own per-cell field, which it reproduces to nine
significant figures. The measured growth is therefore not a property of the interior scheme.
Splitting that field by boundary adjacency locates it: all $484\,344$ interior cells carry a
single value, and the recorded maximum sits on a boundary cell.
\end{abstract}

\section{The question, and why algebra answers it}

METHOD.md section 8.2 records the measurement in Table~\ref{tab:measured}.

\begin{table}[h]
\centering
\caption{The measured transport spectrum on hexahedral meshes (METHOD.md 8.2,
\texttt{blockMesh hex} rows; 2000 power iterations, clip off, frozen velocity).}
\label{tab:measured}
\begin{tabular}{lrrrrr}
\toprule
mesh & cells & $h$ [m] & Co & $\rho(B)$ & $\Lambda_{\max}$ \\
\midrule
\texttt{blockMesh} hex, $n=32$ & 65\,536  & 7.81e-05 & 0.0164 & 1.00439 & 1.0527 \\
\texttt{blockMesh} hex, $n=48$ & 221\,184 & 5.21e-05 & 0.0164 & 1.00441 & ---    \\
\texttt{blockMesh} hex, $n=64$ & 524\,288 & 3.91e-05 & 0.0164 & \textbf{1.00441} & 1.0527 \\
\bottomrule
\end{tabular}
\end{table}

Two features set up this report. First, $\rho-1$ is invariant under refinement at fixed
Courant number: $4.39$e-3, $4.41$e-3, $4.41$e-3 across a two-fold range of $h$ and an
eight-fold range of cell count. Second, the same section records that $\Lambda$ is not a
reliable proxy for $\rho$ --- refining \texttt{pMesh} raises $\Lambda_{\max}$ by $4.3\,\%$
while $\rho-1$ \emph{falls} by $7\,\%$.

On a uniform hexahedral mesh, away from boundaries, the semi-Lagrangian update is a linear
operator with \emph{constant coefficients}: the stencil geometry, the weights and the
displacement are identical at every cell. Such an operator on a periodic uniform mesh is
diagonalised exactly by Fourier modes, so its amplification factor is available in closed
form. The measured number must be reproducible with algebra rather than with 2000 power
iterations on half a million cells.

\section{The operator, as implemented}

The production reconstruction is
\texttt{uncachedQuadraticWeightedLeastSquaresReconstruction}. Four properties fix the
operator completely.

\paragraph{The fit is anchored.} The right-hand side is assembled as the difference
$\psi_i - \psi_c$ over a \emph{constant-free} basis, and \texttt{evaluateRaw} starts from
$\psi_c$, so $R_c(x_c) = \psi_c$ holds by construction. The basis, in the code's order, is
\begin{equation}
  b(d) \;=\; \bigl(\, d_a \;;\; \tfrac{1}{2} d_a^2 \;;\; d_a d_b,\ a<b \,\bigr),
  \qquad
  n_{\text{coeff}} = 2D + \tfrac{1}{2}D(D-1),
\end{equation}
so $n_{\text{coeff}} = 2$ in 1D, $5$ in 2D and $9$ in 3D.

\paragraph{The weights are inverse distance from the cell centre.}
$w_i = 1/\max(|x_i - x_c|,\varepsilon)$, exponent one, with the normal equations using
$w_i^2$. The anchor is the arrival cell centre, not the departure foot point.

\paragraph{The stencil is one ring.} The cell-point-cell stencil gives the full $3^D$ cube
minus the owner: $2$ neighbours in 1D, $8$ in 2D, $26$ in 3D. Physical boundary faces enter
as extra data slots; the interior stencil is untouched by that option.

\paragraph{The update.} With $M = \sum_j w_j^2\, b(d_j)\, b(d_j)^{\!\top}$ and
$d = x_d - x_c$ the departure displacement,
\begin{equation}
  \psi_c^{n+1}
  \;=\; \Bigl(1 - \sum_j g_j\Bigr)\psi_c \;+\; \sum_j g_j \psi_j,
  \qquad
  g_j \;=\; w_j^2\, b(d)^{\!\top} M^{-1} b(d_j).
  \label{eq:update}
\end{equation}

\section{Every $h$ cancels, which derives the measured invariance}
\label{sec:hcancel}

Write $d_j = h s_j$ with integer offsets $s_j \in \{-1,0,1\}^D\setminus\{0\}$ and $d = h c$,
with $c$ the displacement in cell widths. The basis scales as $b(h s) = S_h\, b(s)$ with
$S_h = \operatorname{diag}(h,\dots,h,\,h^2,\dots,h^2)$, and the weights as
$w_j^2 = h^{-2}|s_j|^{-2}$. Hence $M = h^{-2} S_h M_1 S_h$ with
$M_1 = \sum_j |s_j|^{-2} b(s_j) b(s_j)^{\!\top}$, and substituting into \eqref{eq:update}
every factor of $h$ and $S_h$ cancels identically:
\begin{equation}
  \boxed{\;
  g_j(c) \;=\; \frac{b(c)^{\!\top} M_1^{-1} b(s_j)}{|s_j|^{2}} \;}
  \label{eq:gj}
\end{equation}

\derived{} The weights of the update depend \emph{only} on the normalised displacement $c$
and on the integer stencil offsets. They do not depend on $h$ at all. This derives the
measured $h$-invariance of $\rho-1$ at fixed Courant number rather than observing it.

With the cell-centre anchor $M_1$ does not depend on $c$, so it is inverted once and every
$g_j$ is an exact polynomial in $c$: the whole calculation is rational arithmetic.

\section{The symbol}

Inserting $\psi_i = \exp(\mathrm{i}\,\theta\cdot i)$ into \eqref{eq:update},
\begin{equation}
  G(\theta) \;=\; \Bigl(1 - \sum_j g_j\Bigr) \;+\; \sum_j g_j \, e^{\mathrm{i}\,\theta\cdot s_j},
  \qquad
  \rho \;=\; \max_{\theta\in[-\pi,\pi]^D} |G(\theta)| .
  \label{eq:symbol}
\end{equation}
For comparison we also report the Lebesgue constant the code writes as
\texttt{slFitAmplification},
\begin{equation}
  \Lambda \;=\; \Bigl|1 - \sum_j g_j\Bigr| \;+\; \sum_j |g_j| ,
  \label{eq:lambda}
\end{equation}
a row-norm bound on the same weights. $\Lambda$ bounds one step in one cell in the maximum
norm; $\rho$ decides repeated stability in $L^2$. Section~\ref{sec:results} shows how far
apart they are.

\section{Validation}

\subsection{Gate one: 1D must be Lax--Wendroff, symbolically}
\label{sec:validate}

In 1D the stencil is $s\in\{-1,+1\}$ and $M_1$ is $2\times 2$. Putting $c = -C$, so the foot
lies upstream at $x_c - Ch$, equation~\eqref{eq:gj} gives \emph{symbolically}
\begin{equation}
  g_{-1} = \frac{C(C+1)}{2}, \qquad
  g_{+1} = \frac{C(C-1)}{2}, \qquad
  1 - \textstyle\sum_j g_j = 1 - C^2 ,
\end{equation}
which is exactly Lax--Wendroff, and from them
\begin{equation}
  \Lambda = 1 + C - C^2 > 1,
  \qquad
  |G|^2 = 1 - 4C^2\bigl(1-C^2\bigr)\sin^4(\theta/2) \;\le\; 1 .
  \label{eq:lw}
\end{equation}
\derived{} Both reproduce METHOD.md 8.2 item 3 exactly. The command
\texttt{sl\_amplification\_symbol.py validate-1d} checks all five identities symbolically and
exits non-zero on any mismatch. It passes.

Equation~\eqref{eq:lw} is also the warning: $\Lambda > 1$ while $|G| \le 1$. The scheme has
negative weights and is still $L^2$-stable, so a bound on $\Lambda$ cannot establish
instability.

\subsection{Gate two: an independent power iteration}

The symbol is an algebraic construction, so it is checked against a procedure that shares
none of its machinery: the interior stencil is applied repeatedly to a random field on a
periodic $24^3$ grid by \texttt{numpy.roll}, renormalising each step --- the same procedure
as \texttt{leiaTestTransportSpectrum}, restricted to the translation-invariant operator.

\measured{} The iteration never exceeds one. It converges to $1$ from below, giving
$0.99988$ to $0.99992$ after 600 iterations at Courant numbers $0.0164$, $0.25$, $0.5$ and
$0.9$, along $x$ and along $(1,1,1)$. The approach from below is expected: every
non-constant mode is damped, so the iteration is converging onto the constant mode, whose
eigenvalue is exactly one.

\subsection{Gate three: the solver's own per-cell field}
\label{sec:gate3}

The decisive check is against the implementation. The probe
\texttt{workflow/scripts/fit\_amplification\_probe.sh} was run on the \emph{same} mesh as
Table~\ref{tab:measured}, the $128\times 64\times 64$ uniform hexahedral box of
\texttt{popinetTranslating3D} at $n = 64$, with the recorded probe displacement. It
reproduces the recorded maximum, $\Lambda_{\max} = 1.0527251$, and reports no cell demoted to
a linear fit (smallest scaled Cholesky pivot $0.6367$, against the tolerance $0.3$).

The resulting field is split by boundary adjacency in Table~\ref{tab:split}.

\begin{table}[h]
\centering
\caption{The per-cell \texttt{slFitAmplification} field of the $n=64$ hexahedral box, split
into cells that own a boundary face and cells that do not.}
\label{tab:split}
\begin{tabular}{lrrrr}
\toprule
 & cells & min & max & std \\
\midrule
interior          & 484\,344 & 1.01639260 & 1.01639260 & $2.2\times10^{-16}$ \\
boundary-adjacent &  39\,944 & 1.01643120 & \textbf{1.05272510} & $8.1\times10^{-3}$ \\
\bottomrule
\end{tabular}
\end{table}

\measured{} Every one of the $484\,344$ interior cells carries the \emph{same} value, to
machine epsilon --- one distinct value at eight decimal places. That is the
constant-coefficient structure of section~\ref{sec:hcancel}, confirmed by the solver.

\derived{} Evaluating \eqref{eq:gj} at the exact displacement of that run,
$\mathrm{Co} = 6.39947740050255\times10^{-7} / (0.0025/64) = 0.016382662$, gives
$\Lambda = 1.01639260$ against the code's $1.01639260$: a difference of
$2.6\times 10^{-9}$, or nine significant figures. The symbolic model of the operator is the
operator.

\measured{} The global maximum, $1.0527251$, sits on a \textbf{boundary} cell. The
boundary-adjacent set takes eight distinct values; the interior takes one.

\section{Why 1D cannot answer the question}

Table~\ref{tab:ladder} is why the derivation must be carried to three dimensions.

\begin{table}[h]
\centering
\caption{The fit is over-determined only in two and three dimensions.}
\label{tab:ladder}
\begin{tabular}{crrl}
\toprule
$D$ & neighbours & coefficients & character \\
\midrule
1 & 2  & 2 & \textbf{exactly determined}: an interpolant; the weights cancel \\
2 & 8  & 5 & over-determined by 3 \\
3 & 26 & 9 & over-determined by 17 \\
\bottomrule
\end{tabular}
\end{table}

\derived{} In 1D the number of neighbours equals the number of coefficients, so the
least-squares problem has a unique exact solution and the fit \emph{interpolates}. The
weights $w_i$ cancel out of \eqref{eq:gj} entirely --- they may be changed arbitrarily
without changing the scheme. The one-dimensional analysis, including the Lax--Wendroff
reassurance of \eqref{eq:lw}, therefore says nothing about the least-squares character of the
production operator. Over-determination, and with it any effect of the weight choice, is a
two- and three-dimensional phenomenon, and in 3D it is large: 26 data for 9 unknowns.

\section{Results}
\label{sec:results}

\subsection{The interior operator does not amplify, anywhere}

\derived{} Over the ninety-six configurations of
\texttt{\tables/sl\_amplification\_symbol.csv} --- dimensions $1,2,3$; weight anchors
cell-centre and foot-point; displacement along $x$ and along $(1,1,1)$; Courant numbers
$0.0164$, $0.05$, $0.1$, $0.2$, \dots, $0.95$ --- the result is
\begin{equation}
  \rho \;=\; 1.00000000 \qquad\text{in every case, with no exception,}
\end{equation}
while $\Lambda$ ranges from $1.016$ to $2.022$. A $96^3$ grid search plus $4\times10^5$
random samples finds no mode with $|G| > 1$. The maximum of $|G|$ is attained at $\theta = 0$
and equals one, which is the consistency requirement that the scheme reproduces a constant.

\begin{table}[h]
\centering
\caption{Representative rows: the interior operator at the production displacement and at
large Courant number. $\Lambda$ moves a great deal; $\rho$ does not move at all.}
\begin{tabular}{llrrr}
\toprule
$D$ & anchor & Co & $\Lambda$ & $\rho$ \\
\midrule
3 & cell-centre & 0.0164 & 1.016410 & 1.00000000 \\
3 & cell-centre & 0.50   & 1.509259 & 1.00000000 \\
3 & cell-centre & 0.90   & 2.022308 & 1.00000000 \\
3 & foot-point  & 0.0164 & 1.016321 & 1.00000000 \\
3 & foot-point  & 0.50   & 1.357130 & 1.00000000 \\
\midrule
\multicolumn{3}{l}{\emph{measured, whole mesh (Table~\ref{tab:measured})}} & \emph{1.0527} & \emph{1.00441} \\
\bottomrule
\end{tabular}
\end{table}

\paragraph{The neutral modes depend on the flow direction.}
For a displacement along $(1,1,1)$ the constant is the \emph{only} neutral mode: the largest
$|G|$ away from $\theta = 0$ is $0.99999998$, so every non-constant mode is strictly damped.
For a displacement along $x$, 191 of $96^3$ sampled modes are neutral to $10^{-12}$ --- all
of them modes with $\theta_x = 0$, that is, constant along the flow. Such a mode is not
transported and the transverse fit reproduces it exactly. It is never amplified, but it is
also never damped.

\subsection{Every boundary stencil amplifies}

The same calculation is carried out for the stencils that occur at a box boundary, with the
code's rank ladder applied: where the weighted normal matrix is singular the fit is demoted
to linear, as \texttt{quadraticPivotTol} does in the implementation. The quantity reported is
the symbol of the operator that would result if that stencil were applied everywhere. It is a
local amplification indicator, \emph{not} a spectral radius, because a boundary breaks
translation invariance.

\begin{center}
\begin{tabular}{lrrrr}
\toprule
 & \multicolumn{2}{c}{along $x$} & \multicolumn{2}{c}{along $(1,1,1)$} \\
stencil & $\Lambda$ & $\max|G|$ & $\Lambda$ & $\max|G|$ \\
\midrule
interior, 26 neighbours, quadratic               & 1.0164 & \textbf{1.000000} & 1.0174 & \textbf{1.000000} \\
face cell, 17 neighbours, no face slot, LINEAR   & 1.0328 & 1.032800 & 1.0284 & 1.018937 \\
face cell, 17 neighbours $+$ 1 face, quadratic   & 1.0106 & 1.010575 & 1.0190 & 1.006193 \\
edge cell, 11 neighbours $+$ 2 faces, quadratic  & 1.0118 & 1.010575 & 1.0207 & 1.012744 \\
corner cell, 7 neighbours $+$ 3 faces, quadratic & 1.0123 & 1.011028 & 1.0191 & 1.014504 \\
\bottomrule
\end{tabular}
\end{center}

\derived{} The interior cell is the only stencil whose symbol does not exceed one.

\paragraph{A boundary cell cannot support a quadratic fit without its boundary face.}
With the face slot excluded, a face cell's stencil has only two distinct levels in the
wall-normal direction. The monomials $d_x$ and $\tfrac12 d_x^2$ are then collinear on the
data, the normal matrix is exactly singular, and the fit demotes to linear. Adding the
boundary face at $h/2$ restores a third level and with it the quadratic rank. This is
confirmed by the probe of section~\ref{sec:gate3}: with boundary faces included --- the
default --- \emph{no} cell of the $n=64$ box is demoted, and the smallest scaled pivot is
$0.6367$, comfortably above the $0.3$ tolerance.

\derived{} \textbf{Including the boundary faces is the better of the two options, by a factor
of three}: the linear fallback gives $\max|G| = 1.0328$, including the face gives $1.0106$.
The \texttt{zeroGradient} face value is not a sample of the true field, and the reason to
keep it is not accuracy --- it is that the slot carries the rank.

\opn{} The modelled boundary stencils reach $1.033$, while the measured maximum is $1.0527$.
The models above are of a plain box with one boundary condition; the probe case has
\texttt{inlet}, \texttt{outlet} and \texttt{walls} patches, and the boundary-adjacent set
takes eight distinct values. Which patch and which cell attain the maximum is not settled
here.

\section{What this changes}

\begin{enumerate}
\item \derived{}\,\measured{} \textbf{The interior semi-Lagrangian operator on a uniform
  hexahedral mesh does not amplify.} $\rho = 1$ exactly in all ninety-six configurations,
  confirmed symbolically, by an independent power iteration, and against the solver's own
  per-cell field to nine significant figures. The measured $\rho = 1.00441$ is not a
  property of the interior scheme.

\item \measured{} \textbf{The growth is at the boundary.} All $484\,344$ interior cells of
  the production mesh carry one value; the recorded maximum sits on a boundary cell. The
  search for the mechanism should move there.

\item \derived{} \textbf{$\Lambda$ and $\rho$ do not merely differ in accuracy; they differ
  in kind.} At Courant number $0.9$ the interior fit has $\Lambda = 2.022$ and $\rho = 1$
  exactly. A $\Lambda$ of two on an interior hexahedral cell is not evidence of anything.
  This extends the one-dimensional Lax--Wendroff argument of METHOD.md 8.2 item 3 to the
  production three-dimensional operator, where it had not been established.

\item \derived{} \textbf{Moving the weights to the foot point does not change $\rho$.} It
  changes $\Lambda$ substantially --- at Courant $0.5$ along $x$, from $1.509$ to $1.357$ ---
  but $\rho$ is one for both anchors at every Courant number tested. On uniform interior
  hexahedra the weight anchor is not the mechanism, and this is known before any solver code
  is written.

\item \opn{} \textbf{A mesh-aligned displacement leaves a family of undamped modes.} Whether
  they are excited in practice, and by what, is not settled here.
\end{enumerate}

\section{Reproduction}

Every number is produced by one committed script with no case, no mesh and no solver, except
Table~\ref{tab:split}, which needs the probe:

\begin{verbatim}
python3 workflow/scripts/sl_amplification_symbol.py validate-1d     # the gate
python3 workflow/scripts/sl_amplification_symbol.py rho --dim 3 --courant 0.0164 --axis x
python3 workflow/scripts/sl_amplification_symbol.py weights --dim 3 --courant 0.0164
python3 workflow/scripts/sl_amplification_symbol.py sweep           # the curated table
python3 docs/analysis-amplification-sl/sympy/boundary_lambda.py     # the boundary stencils

workflow/scripts/fit_amplification_probe.sh \
    studies/popinet3D_La12000_hex_smoke4/popinetTranslating3D_00000 reuse <out>
\end{verbatim}

The sweep writes \texttt{\tables/sl\_amplification\_symbol.csv}. The validation gate exits
non-zero on failure and must pass before any other number is quoted.

\end{document}
